\section*{अध्याय \ref{ch:g10:drainage-basins} --- अपवाह द्रोणियाँ, नदियाँ और बाढ़}

\begin{solution}{exo:g10:drainage-basins:1}
नदी के किसी बिंदु की \omterm{def:g10:drainage-basins:drainage-basin}{अपवाह द्रोणी} भूमि का वह पूरा क्षेत्र है जिसका
अपवाह उस बिंदु से ऊपर की ओर नदी में बहता है और उससे गुज़रता है। \omterm{def:g10:drainage-basins:drainage-basin}{जल-विभाजक}
वह रेखा है, जो शिखरों के अनुदिश चलती है और दो पड़ोसी द्रोणियों को अलग करती है।
किसी \omterm{def:g10:drainage-basins:drainage-basin}{जल-विभाजक} के दोनों ओर एक मीटर की दूरी पर गिरने वाली दो बूँदें विपरीत
ढलानों पर, दो अलग-अलग द्रोणियों में बहती हैं: उनकी नदियाँ अलग-अलग समुद्रों तक पहुँच सकती हैं।
\end{solution}

\begin{solution}{exo:g10:drainage-basins:2}
\omterm{def:g10:drainage-basins:discharge}{विसर्जन} जल का वह आयतन है जो हर सेकंड नदी की किसी अनुप्रस्थ काट से होकर
गुज़रता है, \unit{m^3/s} में। \qty{1}{m^3/s} हर
सेकंड \qty{1000}{L} है; एक दिन में $1000 \times 86\,400 = 86.4$ मिलियन लीटर।
\end{solution}

\begin{solution}{exo:g10:drainage-basins:3}
बाढ़ से पहले नदी लगभग \qty{330}{m^3/s} ले जा रही थी (लगभग
\num{240} और \qty{390}{m^3/s} के बीच)। शिखर 21 जनवरी को \qty{9231}{m^3/s}
था। चढ़ाव 18 जनवरी (\qty{362}{m^3/s}) को शुरू हुआ: इसमें तीन
दिन लगे।
\end{solution}

\begin{solution}{exo:g10:drainage-basins:4}
$1/50 = 0.02$, \qty{2}{\percent}। नहीं: वर्ष स्वतंत्र हैं; इस वर्ष भी संभावना
\qty{2}{\percent} ही है, पिछले वर्ष चाहे जो हुआ हो।
\end{solution}

\begin{solution}{exo:g10:drainage-basins:5}
एक घन फ़ुट $0.3048^3 \approx \qty{0.02832}{m^3}$ है, इसलिए
$593\,000 \times 0.02832 \approx \qty{16800}{m^3/s}$। अमेज़न
$206\,000/16\,800 \approx 12$ गुना अधिक ले जाती है।
\end{solution}

\begin{solution}{exo:g10:drainage-basins:6}
पट्टियों के क्षेत्रफल: $1.5 \times 0.3 = 0.45$, $0.90$, $1.05$ और \qty{0.60}{m^2}:
$A = \qty{3.0}{m^2}$। \omterm{def:g10:drainage-basins:discharge}{विसर्जन}: $0.2 \times 0.45 = 0.09$, $0.5 \times 0.90
= 0.45$, $0.6 \times 1.05 = 0.63$, $0.3 \times 0.60 = 0.18$: $Q =
\qty{1.35}{m^3/s}$। माध्य वेग $1.35/3.0 = \qty{0.45}{m/s}$।
\end{solution}

\begin{solution}{exo:g10:drainage-basins:7}
$66 \times 10^6 \times 0.020 = \qty{1.32e6}{m^3}$। दो दिन
$2 \times 86\,400 = \qty{172800}{s}$ हैं: औसतन $1.32 \times 10^6/172\,800 \approx
\qty{7.6}{m^3/s}$।
\end{solution}

\begin{solution}{exo:g10:drainage-basins:8}
क्रम से: 120, 88, 65, 59, 51, 47, 42, 38, 35। $n = 9$ के साथ $T = 10/m$:
सबसे बड़े (120) का $T = 10$ वर्ष, दूसरे (88) का 5 वर्ष, तीसरे (65) का
लगभग 3.3 वर्ष। माध्यिका, 51, का क्रम 5 है और $T = 2$ वर्ष: वह दो में से एक वर्ष
बराबर होता है या पार होता है।
\end{solution}

\begin{solution}{exo:g10:drainage-basins:9}
नगर का \omterm{def:g10:drainage-basins:hydrograph}{जलालेख} जल्दी और कहीं अधिक ऊँचा चढ़ता है, और तेज़ी से गिरता है:
तीव्र ढलानों, छतों और सड़कों पर थोड़ा ही जल रिसता है और अपवाह नाले तक
तेज़ी से और एक साथ पहुँचता है; उसका \omterm{def:g10:drainage-basins:hydrograph}{विलंब-काल} छोटा है। वनाच्छादित द्रोणी का
\omterm{def:g10:drainage-basins:hydrograph}{जलालेख} देर से, एक नीचे, चौड़े शिखर तक चढ़ता है: मृदा और वन
वर्षा का बड़ा भाग सोखकर उसे धीरे-धीरे छोड़ते हैं, और हल्की ढलानों पर जल
धीरे बहता है। तूफ़ान के बाद उसका \omterm{def:g10:drainage-basins:hydrograph}{आधार-प्रवाह} अधिक ऊँचा बना रहता है।
\end{solution}

\begin{solution}{exo:g10:drainage-basins:10}
$9231 \times 86\,400 \approx \qty{7.98e8}{m^3}$, जो
$29\,900 \times 10^6 = \qty{2.99e10}{m^2}$ पर फैला है: लगभग \qty{0.027}{m},
यानी \qty{27}{mm}। \omterm{def:g10:drainage-basins:hydrograph}{आधार-प्रवाह} $330 \times 86\,400/2.99 \times 10^{10}
\approx \qty{0.001}{m}$ देता है, यानी लगभग \qty{1}{mm} प्रतिदिन: शिखर-दिवस लगभग
28 गुना अधिक ले गया।
\end{solution}

\begin{solution}{exo:g10:drainage-basins:11}
$P = 1 - 0.9^{10} \approx 1 - 0.349 = 0.651$: लगभग \qty{65}{\percent}। इसके
बिल्कुल न आने की प्रायिकता $0.9^{10} \approx 0.35$ है, यानी लगभग
\qty{35}{\percent}।
\end{solution}

\begin{solution}{exo:g10:drainage-basins:12}
\omterm{def:g10:drainage-basins:hydrograph}{जलालेख} तीन दिन लंबा और \qty{1200}{m^3/s} ऊँचा एक त्रिभुज है: उसका
क्षेत्रफल $\tfrac12 \times 1200 \times 3 \times 86\,400 \approx
\qty{1.56e8}{m^3}$, यानी \num{156} मिलियन घन मीटर, है।
\qty{800}{m^3/s} से ऊपर का भाग इसी आकार का एक छोटा त्रिभुज है, जो ऊँचाई ($400$, $1200$ के बजाय) और अवधि (तीन के बजाय एक दिन) में तीन गुना
छोटा है: $\tfrac12 \times 400 \times 86\,400 \approx \qty{1.73e7}{m^3}$,
यानी \num{17} मिलियन घन मीटर। यह जलाशय द्वारा संचित किए जा सकने वाले \num{60} मिलियन से
कम है: हाँ।
\end{solution}

\begin{solution}{exo:g10:drainage-basins:13}
किनारों तक भरी नदी $800/(100 \times 2) = \qty{4}{m}$ गहरी है।
(1) बाढ़ का मैदान \qty{900}{m} पर \qty{0.5}{m/s} से
\qty{400}{m^3/s} ले जाता है: उसकी गहराई $400/(900 \times 0.5) \approx \qty{0.9}{m}$ है।
(2) तटबंधों के बीच: $1200/(100 \times 2) = \qty{6}{m}$ गहरी, अपने किनारों से \qty{2}{m}
ऊपर, जो तटबंधों के बिना बाढ़ के मैदान की गहराई के दुगुने से भी
अधिक है।
\end{solution}

\begin{solution}{exo:g10:drainage-basins:14}
हमें $(1 - 1/T)^{30} = 0.5$ चाहिए, इसलिए $1 - 1/T = 0.5^{1/30} \approx 0.97716$,
$1/T \approx 0.02284$ और $T \approx 44$ वर्ष।
\end{solution}

\begin{solution}{exo:g10:drainage-basins:15}
हर वर्ग अपना जल अपने सबसे निचले पड़ोसी को भेजता है (पंक्तियाँ उत्तर से दक्षिण,
स्तंभ पश्चिम से पूर्व):
\begin{center}
\begin{tabular}{|c|c|c|c|}
\hline
निकास & पश्चिम & पश्चिम & पश्चिम \\ \hline
उत्तर & पश्चिम & पश्चिम & दक्षिण \\ \hline
उत्तर & उत्तर & दक्षिण & दक्षिण \\ \hline
उत्तर & पूर्व & पूर्व & निकास \\ \hline
\end{tabular}
\end{center}
तीरों का अनुसरण करने पर वर्ग 47, 39, 44, 37, 46 और 30
दक्षिण-पूर्वी निकास में बहते हैं: 6 वर्ग, \qty{0.06}{km^2}। बाक़ी 10 वर्ग
उत्तर-पश्चिमी निकास में बहते हैं: \qty{0.10}{km^2}। \omterm{def:g10:drainage-basins:drainage-basin}{जल-विभाजक} पूर्वी
किनारे से 50 और 47 के बीच से चलता है, फिर 48 और 47 के बीच, 48 और 44 के बीच,
49 और 44 के बीच, 49 और 46 के बीच, और 51 और 46 के बीच से होकर दक्षिणी किनारे तक।
\end{solution}

\begin{solution}{pb:g10:drainage-basins:1}
\textbf{1.} $11\,560 \times 1.609^2 \approx \qty{29900}{km^2}$।

\textbf{2.} केंद्र से ऊपर की ओर की द्रोणी का \omterm{def:g10:drainage-basins:drainage-basin}{जल-विभाजक}।

\textbf{3.} $484\,000 \times 0.3048^3 \approx \qty{13700}{m^3/s}$।

\textbf{4.} $13\,705/3341 \approx 4.1$ गुना।

\textbf{5.} $A = 400 \times 8 = \qty{3200}{m^2}$ और $v = Q/A = 13\,705/3200
\approx \qty{4.3}{m/s}$।

\textbf{6.} लगभग \qty{330}{m^3/s}।

\textbf{7.} \qty{9231}{m^3/s}, 21 जनवरी को।

\textbf{8.} तीन दिन।

\textbf{9.} $9231 \times 86\,400/(29\,900 \times 10^6) \approx
\qty{0.027}{m}$: लगभग \qty{27}{mm}।

\textbf{10.} नदी तक नया जल पहुँचा: द्रोणी पर एक और वर्षा
(या पिघलती बर्फ़) ने अपवाह की एक दूसरी, छोटी लहर भेजी।

\textbf{11.} शिखर के सत्रह दिन बाद। चढ़ाव नालों के पास की ढलानों के
अपवाह से आता है, जो जल्दी और एक साथ पहुँचता है; अवरोही भुजा पर
नदी को अब भी द्रोणी के सबसे दूर के भागों का जल और मृदा से धीरे-धीरे
निकलता जल मिलता रहता है।

\textbf{12.} $T = n + 1 = 96$ वर्ष।

\textbf{13.} $96/2 = 48$ वर्ष।

\textbf{14.} $96/3 = 32$ वर्ष।

\textbf{15.} $96/4 = 24$ वर्ष और $96/5 = 19.2$ वर्ष। बहुत नहीं: दोनों ही स्थितियों में
यह \omterm{def:g10:drainage-basins:discharge}{विसर्जन} लगभग बीस वर्षों में एक बार लौटता है।

\textbf{16.} पाँच वर्ष (1936, 1937, 1943, 1972 और 1996): $5/95 \approx
0.053$, यानी प्रति वर्ष लगभग \qty{5}{\percent}, जो लगभग
$96/5 \approx 19$ वर्ष का \omterm{def:g10:drainage-basins:return-period}{पुनरावृत्ति काल} है।

\textbf{17.} $1/100 = 0.01$, \qty{1}{\percent}।

\textbf{18.} $1 - 0.01 = 0.99$।

\textbf{19.} $0.99^{30} \approx 0.740$।

\textbf{20.} $1 - 0.740 = 0.26$: \textbf{30 वर्षों में कम से कम एक
100-वर्षीय बाढ़ की संभावना लगभग \qty{26}{\percent} है}, यानी चार में एक।
\end{solution}
